% Generated draft: requires mathematical and source review.
\documentclass[11pt]{article}
\usepackage[margin=25mm]{geometry}
\usepackage{fontspec}
\setmainfont{Latin Modern Roman}
\usepackage{amsmath,amssymb,mathtools,unicode-math}
\setmathfont{Latin Modern Math}
\usepackage{enumitem,needspace,xurl,hyperref}
\setlength{\parindent}{0pt}
\setlength{\parskip}{5pt}
\setlength{\emergencystretch}{3em}
\newcommand{\N}{\mathbb{N}}
\newcommand{\Z}{\mathbb{Z}}
\newcommand{\Q}{\mathbb{Q}}
\newcommand{\R}{\mathbb{R}}
\newcommand{\C}{\mathbb{C}}
\newcommand{\F}{\mathbb{F}}
\newcommand{\A}{\mathbb{A}}
\newcommand{\PP}{\mathbb{P}}
\newcommand{\E}{\mathbb{E}}
\newcommand{\eps}{\varepsilon}
\newcommand{\dd}{\,\mathrm{d}}
\begin{document}
\section*{UnsolvedMath ID: 30004952}
\textbf{Source problem code:} catalogue entry
\subsection{Definitions and mathematical statement}
Let \(n\) denote the ambient Euclidean dimension and let \(S^{n-1}=\{u\in\R^n:|u|=1\}\). The supplied record does not specify a range for \(n\). We interpret the extracted notation \(K\in\mathcal K_o^n\) in its usual sense: \(K\subset\R^n\) is a compact convex set with nonempty interior and \(0\in\operatorname{int}K\). This interpretation of the damaged source notation requires confirmation against the report.

For such a body, define its radial function by
\[
\rho_K(u)=\max\{t\geq0:tu\in K\},\qquad u\in S^{n-1}.
\]
For clarity, use the following conventional definition of Aleksandrov integral curvature. Let \(\sigma\) be ordinary, unnormalized spherical surface measure. For \(\sigma\)-almost every \(v\in S^{n-1}\), the supporting hyperplane with outer normal \(v\) meets \(K\) at a unique point \(x_K(v)\). Put \(r_K(v)=x_K(v)/|x_K(v)|\), with arbitrary values on the exceptional null set, and define
\[
J(K,A)=\int_{S^{n-1}}\mathbf 1_A(r_K(v))\,\mathrm d\sigma(v)
\]
for every Borel set \(A\subseteq S^{n-1}\). The record names \(J\) but does not itself supply this construction or its normalization; these conventions require source review.

For each \(p\in\R\), define the \(L_p\) integral curvature measure by the supplied identity
\[
\mathrm dJ_p(K,\cdot)=\rho_K^p\,\mathrm dJ(K,\cdot),
\qquad
J_p(K,A)=\int_A\rho_K(u)^p\,\mathrm dJ(K,u).
\]
In particular, \(J_0(K,\cdot)=J(K,\cdot)\).

\textbf{Problem.} For each real \(p\), characterize the finite nonnegative Borel measures \(\mu\) on \(S^{n-1}\) for which there exists \(K\in\mathcal K_o^n\) satisfying
\[
\mu(A)=J_p(K,A)
\quad\text{for every Borel set }A\subseteq S^{n-1}.
\]
The requested characterization must give necessary and sufficient conditions for existence. No symmetry or absolute continuity hypothesis is imposed in the supplied general statement.

The original excerpt also asserts a Monge--Amp\`ere-type formulation when \(\mu\) has a density \(f\). Its displayed equation is corrupted, and the supplied text does not specify the reference measure for \(f\), the role of \(h\), or the required regularity. Consequently, that equation cannot be stated precisely from this record and is not reconstructed here. The background mentions unresolved uniqueness questions, but the stated target above concerns existence; no precise uniqueness question is supplied.
\subsection{Short English statement}
For every real parameter \(p\), determine exactly which finite Borel measures on the sphere are the \(L_p\) integral curvature measures of convex bodies containing the origin in their interior. The origin condition and curvature normalization use conventional interpretations requiring source confirmation. The supplied density-based differential equation is too corrupted to reproduce reliably. This is a statement-only draft requiring human mathematical review.
\subsection{Sources}
\begin{itemize}
\item Supplied local UnsolvedMath registry record: permanent ID \(30004952\), exact source problem code \texttt{catalogue entry}, title ``The \(L_p\) Aleksandrov Problem for Integral Curvature Measures.'' Collection: Oberwolfach Reports Open Problems. The record gives proposed year 2021 and names no individual proposer.
\item Supplied authoritative report reference: \emph{Convex Geometry and its Applications} (2022), Oberwolfach Reports, pp.~3187--3266, DOI \texttt{10.4171/owr/2021/59}. \url{https://doi.org/10.4171/owr/2021/59}. The record attributes the problem to workshop open-problem discussions collected in this report.
\item Huang--Lutwak--Yang--Zhang, supplied literature reference: \url{https://doi.org/10.4310/jdg/1536285625}.
\item Zhao on origin-symmetric polytopes, supplied literature reference: \url{https://doi.org/10.1090/proc/14568}.
\item Zhao current negative-\(p\) results, supplied literature reference: \url{https://www.yimingzhao.net/files/LpAleksandrov.pdf}.
\item Literature assessment dated exactly \texttt{2026-08-22}: the supplied curation labels the problem partially solved and reports a complete existence characterization for positive \(p\), substantial even-measure results for negative \(p\), including discrete polytopes, and remaining questions for unrestricted negative \(p\) and uniqueness. This is supplied dated curation derived from the OWR record, not an independently verified literature assessment; its time-sensitive status and cited claims require independent verification.
\item The registry labels its statement \texttt{corrected\_verified} and reports checking it against the official Oberwolfach report. This is supplied verification metadata; no source visit or independent mathematical verification is asserted here.
\end{itemize}
\end{document}
